\chapter{The Heap: \texttt{malloc}/\texttt{free} in Ada, the TLSF Way}
\label{ch:heap}

Two of the runtime profiles give you a real dynamic heap, and that heap is now
written entirely in Ada. Both Ada's \texttt{new} and the C library's
\texttt{malloc} land in the same allocator, which is a \emph{Two-Level
Segregated Fit} (TLSF)\footnote{TLSF is the algorithm of M.~Masmano, I.~Ripoll,
A.~Crespo and J.~Real, ``TLSF: A New Dynamic Memory Allocator for Real-Time
Systems,'' \emph{16th Euromicro Conference on Real-Time Systems (ECRTS)}, 2004
(\url{http://www.gii.upv.es/tlsf/}). The implementation here is original Ada; the
design is theirs.} implementation: O(1) allocate, O(1) free-and-coalesce,
and bounded fragmentation. This chapter is the current truth about the heap. The
``first-fit'' descriptions in \S\ref{ch:profiles} and the appendix
(\S\ref{ch:appendix}) describe the allocator that TLSF replaced; what follows
supersedes them.

\section{Why a heap, and what lands here}\index{heap}\index{dynamic allocation}\index{secondary stack}

Ada programs reach the heap by two doors. The first is the language itself: every
\texttt{new T} (and every access-type allocator, and the storage pool behind
unconstrained returns when no other pool is named) ultimately calls the runtime's
allocation primitive, which is wired to \texttt{\_\_gnat\_malloc} / the C
\texttt{malloc} symbol. The second door is C code linked into the image (a
vendor blob, a bit of glue, anything that calls \texttt{malloc}, \texttt{calloc},
\texttt{realloc} or \texttt{free}). Both doors open onto the single allocator
described here.

What does \emph{not} come here is the secondary stack. Ada's secondary stack,
which carries function results whose size is not known to the caller, such as an
unconstrained \texttt{String} return, is a separate, contiguously-managed
region with its own mark/release discipline; it never touches \texttt{malloc}.
Keeping the two apart is what lets the leanest profile run with no heap at all.

That leanest profile is \texttt{light-tasking}: it is configured
\texttt{No\_Secondary\_Stack} and heap-less, with only a bump allocator that can
hand out memory but never reclaim it (see \S\ref{ch:profiles}). A program built
\texttt{light-tasking} that calls \texttt{new} in a loop will eventually run the
bump pointer off the end of the arena, by design because that profile is for
code that is statically sized and never frees. The \texttt{embedded} and
\texttt{full} profiles are the ones that link the real allocator, so on those
\texttt{free} actually returns memory to the pool and a long-running program can
allocate and release indefinitely.

\section{Why first-fit had to go}\index{first-fit allocator}\index{fragmentation}

The allocator TLSF replaced kept a single linked list of free blocks. To satisfy
a request it walked that list from the head looking for the first block large
enough: O(n) in the number of free blocks. Worse, to keep fragmentation in
check it coalesced on every \texttt{free} by scanning the whole list to find the
freed block's physical neighbours, again O(n). For a small internal-RAM heap
with a handful of live blocks, nobody notices. The problem is the motivating case
for this whole chapter: a \emph{large, slow} heap.

The ESP32-S3 can map up to 8\,MiB of external PSRAM, and that PSRAM is an order of
magnitude slower than on-chip SRAM: it runs as DDR with a wait state, so every
header you touch is an off-chip burst. Put a heap there (the
\texttt{HEAP\_PSRAM=1} mode of the on-target test does exactly this) and a list
that may hold thousands of free blocks turns each allocate \emph{and} each free
into a long off-chip walk. The cost is not just big; it is unpredictable, which is
exactly what you do not want on a real-time runtime. TLSF makes both operations
O(1) with a small, fixed number of memory touches, so the heap behaves the same
whether it lives in fast SRAM or slow PSRAM.

\cnote{ESP-IDF answers the slow-vs-fast question with \texttt{heap\_caps\_malloc}
and a set of capability flags (\texttt{MALLOC\_CAP\_INTERNAL},
\texttt{MALLOC\_CAP\_SPIRAM}, \ldots) routing to several heaps. Here there is one
allocator and one arena; \emph{which} arena (leftover DRAM or PSRAM) is chosen
once, at link time, in \texttt{bare\_build.sh}.}

\section{TLSF, taught from the code}\index{TLSF allocator}

The portable allocator lives in \texttt{tlsf\_core.ads}/\texttt{.adb}. It is
deliberately target-free: no lock, no assembly, no linker symbols. It takes an
arena from its caller and manages it. That purity is what lets the identical code
run on the host (\S\ref{sec:heap-validation}). Here is the whole interface:

\begin{lstlisting}
procedure Init (Base : System.Address; Size : Storage_Count);
function  Allocate (N : Storage_Count) return System.Address;
procedure Deallocate (P : System.Address);
function  Reallocate (P : System.Address; N : Storage_Count)
                      return System.Address;
function  Ready return Boolean;
function  Invariants_Hold return Boolean;   --  test-only
\end{lstlisting}

\subsection{Boundary tags: every block knows its neighbours}\index{boundary tags}

The arena is one contiguous chain of physically-adjacent blocks. Each block
carries a header before its payload:

\begin{lstlisting}
type Hdr is record
   Prev_Phys : System.Address;   --  physical previous block (Null = first)
   Size      : Storage_Count;    --  payload bytes (16-aligned)
   Free      : Boolean;
end record;
\end{lstlisting}

Because a block knows both its own \texttt{Size} and the address of its physical
predecessor (\texttt{Prev\_Phys}), the runtime can step forward to the next block
(\texttt{Next\_Phys = B + Hdr\_Sz + Size}) and back to the previous one in
constant time, without any list walk. That is the \emph{boundary tag} idea, and
it is what makes \texttt{free} O(1): on a free we look only at the block itself
and its two immediate physical neighbours, coalescing with whichever of them is
also free.

A header-size subtlety worth flagging because it bit this code: \texttt{Hdr\_Sz}
is a compile-time \emph{static} constant, not derived from
\texttt{Hdr'Object\_Size}. An object-size-based constant would be \emph{elaborated}
into \texttt{.bss}, and the ZFP boot has no \texttt{adainit} to run that
elaboration, so the constant would read as garbage and produce wild pointers.
For the same reason all the mutable state (\texttt{Heads}, the bitmaps,
\texttt{Pool\_Lo}, \texttt{Sentinel}) has zero initializers, which places it in
\texttt{.bss} (zeroed by \texttt{start.S}) rather than \texttt{.data}.

\subsection{Two-level segregated free lists}\index{segregated free lists}

Free blocks are bucketed by size into a two-dimensional grid of lists. The
\emph{first level} (FL) is a power-of-two class: roughly, $\mathrm{FL} =
\lfloor\log_2 \mathrm{size}\rfloor$. The \emph{second level} (SL) subdivides each
power-of-two range linearly into \texttt{SL\_Count} equal sub-ranges. Here that is
$2^5 = 32$ second-level lists per first-level class, across 25 first-level
classes:

\begin{lstlisting}
SL_Log2  : constant := 5;
SL_Count : constant := 2 ** SL_Log2;   --  32 second-level lists
FL_Count : constant := 25;             --  first-level classes
Heads : array (0 .. FL_Count - 1, 0 .. SL_Count - 1) of System.Address;
\end{lstlisting}

\texttt{Heads (FL, SL)} is the head of a doubly-linked list of free blocks whose
sizes all fall in that one bucket. A free block stores its list links in the first
two words of its \emph{payload} (it is free, so the payload is unused),
\texttt{Next\_Free} / \texttt{Prev\_Free}, which is why the minimum payload is
two words wide.

The \texttt{Mapping} function turns a size into its \texttt{(FL, SL)} indices.
Above the small-block threshold it finds the high bit with a find-last-set
(\texttt{FLS}), uses that for FL, and takes the next \texttt{SL\_Log2} bits below
it for SL:

\begin{lstlisting}
F  : constant Integer := FLS (S);                         --  high bit
SL := Integer (Shift_Right (S, F - SL_Log2) and (SL_Count - 1));
FL := F - FL_Shift + 1;
\end{lstlisting}

A worked example with \texttt{FL\_Shift = 9} (which is \texttt{SL\_Log2 + 4}): a
request rounding to 1536 bytes is \texttt{0b110\_0000\_0000}, whose top bit
(\texttt{FLS}) is at position 10. \texttt{SL} takes the five bits starting just
below that bit: $(1536 \gg (10-5)) \mathbin{\mathrm{and}} 31 = 48
\mathbin{\mathrm{and}} 31 = 16$, and $\mathrm{FL} = 10-9+1 = 2$. So 1536-byte
blocks live in bucket \texttt{Heads (2, 16)}. A 2048-byte block
(\texttt{0b1000\_0000\_0000}, top bit 11) lands in \texttt{Heads (3, 0)}, a
different first-level class, exactly as you want at a power-of-two boundary.

\subsection{Bitmaps + find-first-set: O(1) bucket lookup}

The grid would still need scanning if we had to search it. We do not because a
pair of bitmaps records which buckets are non-empty:

\begin{lstlisting}
FL_Bitmap : Unsigned_32;                       --  bit FL set if class FL has any free block
SL_Bitmap : array (0 .. FL_Count - 1) of Unsigned_32;  --  per-class SL occupancy
\end{lstlisting}

To allocate, \texttt{Mapping\_Search} first rounds the request \emph{up} so that
\emph{any} block in the chosen bucket is guaranteed big enough (this is what makes
it ``good-fit'' rather than risking a too-small first block). Then
\texttt{Find\_Suitable} masks off the buckets at or below the requested class and
uses find-first-set (\texttt{FFS}) to jump straight to the smallest non-empty
bucket that fits: first within the same FL class, then, if that class is
exhausted, by consulting \texttt{FL\_Bitmap} for the next non-empty class up. Two
masked bit-scans, no list walk. That is the O(1) \texttt{malloc}.

\texttt{FFS} and \texttt{FLS} here are written as bounded shift loops rather than
a hardware count-leading-zeros instruction. The loops run at most 32 (or 64)
iterations regardless of heap size, so they are O(1) in the heap and keep the core
free of target intrinsics.

\begin{figure}[ht]
\centering
\resizebox{0.92\linewidth}{!}{%
\begin{tikzpicture}[node distance=6mm]
\node[blk] (req) {request\\\scriptsize round up};
\node[blk, right=of req] (map) {\texttt{Mapping}\\\scriptsize size $\to$ (FL,SL)};
\node[blk, right=of map] (ffs) {bitmaps + FFS\\\scriptsize smallest fitting bucket};
\node[blk, right=of ffs] (pop) {pop head\\\scriptsize of \texttt{Heads(FL,SL)}};
\node[blk, right=of pop] (spl) {split?\\\scriptsize remainder $\to$ free};
\draw[flow] (req)--(map); \draw[flow] (map)--(ffs);
\draw[flow] (ffs)--(pop); \draw[flow] (pop)--(spl);
\end{tikzpicture}}
\caption{One \texttt{Allocate}: map the size, find the bucket by bit-scan, take
the head block, and split off any large remainder back into a free bucket. No
step depends on the number of free blocks.}
\label{fig:tlsf-alloc}
\end{figure}

\subsection{Allocate, free, split, coalesce}\index{coalescing}

\texttt{Allocate} rounds the request to a 16-byte multiple (and up to the
two-word minimum payload), finds a bucket, pops the head block off it, and,
if the block is substantially larger than needed, \texttt{Split}s it: the
front becomes the returned block and the tail is re-inserted as a smaller free
block (updating both its and its successor's \texttt{Prev\_Phys}). The leftover
is dropped back into whatever bucket its new size maps to.

\texttt{Deallocate} is the mirror image and the place boundary tags earn their
keep. It marks the block free, then coalesces \emph{forward} (if
\texttt{Next\_Phys} is free, remove it from its bucket and absorb it) and
\emph{backward} (if \texttt{Prev\_Phys} is free, absorb into it). The terminating
\texttt{Sentinel} block (a zero-size, permanently-\emph{used} block at the top
of the arena) is what makes the forward check safe without a bounds test: the
real top block always has a non-free successor to stop at. After coalescing, the
single resulting block is inserted into its bucket and the relevant bitmap bits
are set. Every one of these steps touches a fixed, small number of blocks.

\texttt{Reallocate} shrinks in place when the existing block is already big
enough, and otherwise allocates-copies-frees. \texttt{Init} lays the arena out as
exactly one big free block plus the sentinel.

\section{The target glue: \texttt{bare\_heap}}\index{bare\_heap}\index{malloc}\index{free}

\texttt{tlsf\_core} has no idea it is on a microcontroller. Everything
target-specific lives in \texttt{bare\_heap.ads}/\texttt{.adb}, which is the unit
that actually exports the C ABI. It replaces the old \texttt{bare\_heap.c}.

\begin{lstlisting}
function  Malloc  (N : Interfaces.C.size_t) return System.Address;
pragma Export (C, Malloc, "malloc");

procedure Free    (P : System.Address);
pragma Export (C, Free, "free");

function  Realloc (P : System.Address; N : Interfaces.C.size_t)
                   return System.Address;
pragma Export (C, Realloc, "realloc");

function  Calloc  (Nmemb, Size : Interfaces.C.size_t) return System.Address;
pragma Export (C, Calloc, "calloc");

procedure Task_Stack_Free (Stack, Thread : System.Address);
pragma Export (C, Task_Stack_Free, "__gnat_task_stack_free");
\end{lstlisting}

\subsection{The critical section}\index{critical section}

\texttt{tlsf\_core} is not internally locked, by design, so the host harness
can drive it single-threaded. On target, multiple agents reach the heap: the
environment task, finalizers running RAII cleanup, and any task calling
\texttt{new}. \texttt{bare\_heap} serialises them with a global critical section
built from the Xtensa \texttt{rsil} instruction at level 15 (raise interrupt
level to the maximum, returning the old PS) around each operation, restored with
\texttt{wsr.ps}\,/\,\texttt{rsync}:

\begin{lstlisting}
function Enter_Crit return Unsigned_32 is
   Ps : Unsigned_32;
begin
   Asm ("rsil %0, 15",
        Outputs  => Unsigned_32'Asm_Output ("=r", Ps),
        Volatile => True, Clobber => "memory");
   return Ps;
end Enter_Crit;
\end{lstlisting}

This matches what the old C heap did. \texttt{Calloc} is careful to zero the
returned block \emph{outside} the lock. It only needs the lock to allocate, not
to memset, which keeps the critical section short.

\subsection{An arena chosen at link time}

\texttt{bare\_heap} never names an address. It imports two linker symbols whose
\emph{addresses} are the arena bounds:

\begin{lstlisting}
Heap_Base : Storage_Element
  with Import, Convention => C, External_Name => "__bare_heap_base";
Heap_End  : Storage_Element
  with Import, Convention => C, External_Name => "__bare_heap_end";
\end{lstlisting}

\texttt{bare\_build.sh} supplies these via \texttt{-{}-defsym}, pointing them
either at the leftover internal DRAM after \texttt{.bss}, or (for the large
heap) at the PSRAM region. The Ada is identical either way; only the two
\texttt{-{}-defsym} values differ. The first \texttt{malloc} lazily calls
\texttt{T.Init} over whatever region those symbols delimit (a \texttt{Started}
flag, also \texttt{.bss}-resident, gates this).

\subsection{Freeing a task's own stack}\index{task stack reclamation}

The full profile allocates each task's primary stack on the heap, which raises a
nasty question: when a task terminates, who frees the stack it was \emph{running
on}, and how do we avoid pulling it out from under a context that has not yet
switched away, possibly on the other core? \texttt{Task\_Stack\_Free}
(exported as \texttt{\_\_gnat\_task\_stack\_free}) handles this with a bounded
spin against the runtime's per-core running-thread table:

\begin{lstlisting}
Running : array (0 .. 1) of System.Address
  with Import, Convention => C, External_Name => "__gnat_running_thread_table";

procedure Task_Stack_Free (Stack, Thread : System.Address) is
begin
   for I in 1 .. 2_000_000 loop
      if Running (0) /= Thread and then Running (1) /= Thread then
         Free (Stack);
         return;
      end if;
   end loop;
   --  timed out -- leak rather than risk a use-after-free
end Task_Stack_Free;
\end{lstlisting}

It frees only once the terminating thread is no longer the running thread on
either core. If the thread somehow never clears (the spin bound expires), it
chooses to \emph{leak} the stack rather than risk a use-after-free: a deliberate
``leak-on-timeout'' policy.

\section{Validation}
\label{sec:heap-validation}\index{heap validation}

\texttt{tlsf\_core} is tested off-chip as ordinary Ada because it is
target-free.
\texttt{heap\_stress.adb} is a generic that takes an allocator's six operations
and beats on them: an initial sanity pass (alloc/align/free, OOM returns null,
realloc-grow preserves data), then 300\,000 randomised alloc/free steps over a
64\,KiB arena, each allocation stamped with a unique byte pattern that is
re-verified before every free so any overlap or stale pointer is caught. After a
full drain it asserts the arena coalesced back to (nearly) one block: proof the
free path actually merges. \texttt{allocator\_test.adb} instantiates that harness
over \emph{both} the old first-fit core and \texttt{Tlsf\_Core}, so a behavioural
divergence between them would show up.

The structural check throughout is \texttt{Invariants\_Hold}, run after every
operation. It walks the physical chain from \texttt{Pool\_Lo} to the
\texttt{Sentinel} and verifies: each block is 16-aligned, each
\texttt{Prev\_Phys} points back to the actual previous block, sizes are at least
the minimum and stay within the arena, and (the key correctness property) no
two physically-adjacent blocks are both free (i.e.\ coalescing is complete). The
sentinel must be zero-size, used, and correctly back-linked.

On hardware, \texttt{examples/esp32s3\_heaptest} drives the live
\texttt{malloc}/\texttt{free} symbols directly over the real arena: 50\,000 random
operations, 32 live blocks, the same pattern-stamping idea, reporting
\texttt{PASS}/\texttt{FAIL} over \texttt{ESP32S3.Log}. Its \texttt{HEAP\_PSRAM=1}
build relocates the arena into the 8\,MiB PSRAM at \texttt{0x3D000000}: the
large, slow heap that motivated TLSF in the first place. See \S\ref{ch:hal} for
how PSRAM is brought up before the heap can use it.

\section{Which profiles get it}

\texttt{light-tasking} links neither \texttt{bare\_heap} nor \texttt{tlsf\_core};
it is heap-less and bump-only (\S\ref{ch:profiles}). \texttt{embedded} and
\texttt{full} link both, so \texttt{malloc}/\texttt{free}/\texttt{realloc}/%
\texttt{calloc} (and therefore Ada's \texttt{new}) are backed by TLSF; the
full profile additionally uses \texttt{Task\_Stack\_Free} for per-task primary-%
stack reclamation. And because the same \texttt{tlsf\_core} compiles and runs
host-native, the allocator the chip runs is the very allocator the host test
exercises, not a model of it.
